% !TeX root = project-report-cn.tex
% !TeX program = lualatex
% Build: latexmk -lualatex -interaction=nonstopmode -halt-on-error -outdir=.build/project-report-cn project-report-cn.tex
\documentclass[11pt,a4paper,fontset=fandol,scheme=plain]{ctexart}
\usepackage[margin=21mm,headheight=14pt]{geometry}
\usepackage{lmodern,amsmath,amssymb,graphicx,booktabs,array,tabularx}
\usepackage{xcolor,enumitem,fancyhdr,caption,xurl,placeins,float,needspace}
\usepackage[colorlinks=true,linkcolor=teal,citecolor=teal,urlcolor=teal]{hyperref}
\definecolor{reportblue}{HTML}{18536B}
\hypersetup{pdftitle={不同特征组对四轻子模拟样本筛选效果的影响},pdfauthor={HiggsML 项目}}
\pagestyle{fancy}
\fancyhf{}
\fancyhead[L]{\small\textcolor{reportblue}{HiggsML | 项目报告}}
\fancyhead[R]{\small 2026 年 11 月 8 日}
\fancyfoot[C]{\thepage}
\setlength{\parindent}{0pt}
\setlength{\parskip}{6pt}
\setlength{\emergencystretch}{2em}
\setlist[itemize]{leftmargin=1.4em,itemsep=3pt,topsep=4pt}
\captionsetup{font=small,labelfont=bf,justification=raggedright,singlelinecheck=false}
\captionsetup[table]{name=表}
\captionsetup[figure]{name=图}
\newcolumntype{Y}{>{\raggedright\arraybackslash}X}
\ctexset{
  section={format=\Large\bfseries\color{reportblue},beforeskip=16pt,afterskip=8pt},
  subsection={format=\large\bfseries\color{reportblue},beforeskip=12pt,afterskip=5pt}}
\setcounter{tocdepth}{1}
\renewcommand{\contentsname}{阅读目录}
\raggedbottom
\widowpenalty=10000
\setlength{\textfloatsep}{12pt}
\setlength{\intextsep}{12pt}
\newcommand{\gev}{\,\mathrm{GeV}}

\begin{document}
{\LARGE\bfseries\color{reportblue} 不同特征组对四轻子模拟样本筛选效果的影响\par}
\vspace{6pt}
项目报告 · 基础指标比较版 · 2026 年 11 月 8 日

\section{研究目的与比较范围}

本报告比较使用不同特征组组合训练的网络，考察信号产额保留率、背景产额保留率以及筛选后的信背比怎样变化。
信背比指信号产额与背景产额之比；质量谱则用来检查筛选前后各质量区域的产额分布。

网络输入分为四组：A 为四个轻子的横动量与相对束流的方向，B 为两个轻子对的质量及各对内部的角分离，
C 为四轻子系统的横动量及两个轻子对之间的方位角差，D 为描述产生与衰变方向的五个角度变量。
各组变量可能相关。例如，B 组的轻子对质量由单轻子的四动量计算，因此它与 A 组的单轻子动量可能提供相关的信号、背景区分线索。

在这四组变量的基础上，比较五种输入方案：ABCD 使用全部四组，BCD 去掉 A，ACD 去掉 B，ABD 去掉 C，ABC 去掉 D。
每种方案使用随机种子 42、43、44 分别训练。所有方案使用相同的训练、验证和最终评估样本：
训练样本用于学习网络参数，验证样本用于选择模型检查点和确定分数阈值，最终评估样本用于应用冻结的选择并计算结果。
网络隐藏层结构和阈值确定规则也保持相同。

模型输出后的比较使用加权求和、比例、差值和中位数。
先介绍样本与网络方法，再由两个模型的计算示例过渡到五种方案的完整比较。
正文采用本机新运行 draft001 的受控 MC 结果；来源及核验范围见附录~\ref{app:provenance}。
研究限于教育与技术用途，结论的适用范围见第~\ref{sec:conclusion}~节。

\section{模拟样本与事件选择}

\subsection{信号、背景与研究总体}
研究使用 ATLAS 2020 教育用途 13 TeV 集合中的两份蒙特卡洛（MC）模拟样本\cite{OpenData}。
一份是信号：样本编号 DSID 345060（\texttt{ggH125\_ZZ4lep}），模拟胶子融合产生质量设为 $125\gev$ 的希格斯，
再经 $H\to ZZ^*\to4\ell$ 衰变。另一份是背景：DSID 363490（\texttt{llll}），同样包含四轻子事件。
两份样本的预期产额均换算到积分亮度 $10\,\mathrm{fb}^{-1}$；这里的产额是给定亮度下的预期事件数，
与文件中有多少条模拟记录不同。

使用网络之前，先按相同的物理条件筛选信号和背景样本。
依次检查事件触发和轻子品质，要求恰好四个合格轻子，再构造同味异号轻子对并施加动量、接受度与轻子对质量条件。
本项目只分析含两个电子和两个缪子（$2e2\mu$）且满足 $105\leq m_{4\ell}<140$ GeV 的事件。
电子对和缪子对组成两个候选；质量更接近 $91.1876$ GeV 的一对称为 $Z_1$，另一对为 $Z_2$。
完整条件、执行顺序和文献对应见附录~\ref{app:cuts}，四动量与角度定义见附录~\ref{app:fields}。

“网络筛选前”已经通过这些共同条件；“网络筛选后”仅额外要求通过固定网络的高分阈值。
因此，两者是同一评估总体与其中的子集，不是独立产生的两批事件。

\subsection{样本划分与各子集的职责}

筛选后的模拟事件按 60\%、20\%、20\% 的比例分为训练、验证和最终评估三个集合，分别用于学习参数、选定模型与阈值、计算最终结果。
同一个事件只分入一个集合。
所有输入方案和种子使用相同的划分，因此比较时不会因更换输入而更换评估事件。
验证时使用的 AUC 衡量信号是否通常排在背景之前，后文再说明用它选择检查点的具体规则。

\begin{center}
\small
\begin{tabularx}{\linewidth}{@{}p{.14\linewidth}p{.18\linewidth}Y@{}}
\toprule
子集 & 分配概率 & 在本研究中负责什么 \\
\midrule
训练 & 60\% & 拟合输入标准化参数与网络参数。 \\
验证 & 20\% & 用 AUC 选择检查点，再用背景分数确定阈值。 \\
最终评估 & 20\% & 在模型与阈值冻结后，计算筛选产额及质量谱。 \\
\bottomrule
\end{tabularx}
\end{center}
训练结束后保存的一组网络参数称为检查点。验证样本负责选定检查点和分数阈值，
最终评估样本只应用已经固定的选择；不能看过评估结果后再改模型或阈值。
\begin{table}[H]
\caption{下载的 MC 样本文件与最终评估样本的数量。}
\label{tab:population}
\centering
\small
\begin{tabularx}{\linewidth}{@{}Y p{.23\linewidth} rrr@{}}
\toprule
样本范围 & 计数方式 & 信号 & 背景 & 合计 \\
\midrule
下载的 MC 样本文件 & 文件中的记录条数 & 164,716 & 554,279 & 718,995 \\
最终评估样本（网络筛选前） & 同一个事件只计一次 & 9,049 & 698 & 9,747 \\
\bottomrule
\end{tabularx}
\end{table}
表~\ref{tab:population} 分别列出下载文件中的记录数，以及通过物理筛选并分入最终评估集合的事件数。

\Needspace{13\baselineskip}
\subsection{事件数、物理权重与预期产额}

\label{sec:weights}
\textbf{事件数}是模拟样本中的事件数量，每个物理事件只计一次。
\textbf{预期产额}是模拟所预测的、在 $10\,\mathrm{fb}^{-1}$ 积分亮度下能有多少个事件。

每个模拟事件的\textbf{物理权重}表示它对预期产额的贡献。计算分为以下三步，信号和背景分别进行。

\textbf{第一步：计算每个事件的物理权重。}
从文件读取生成器权重 $w_i^{\rm gen}=\texttt{mcWeight}$，再按样本参数换算：
\[
 w_i^{\rm phys}=\frac{\mathcal L\,\sigma_{\rm eff}}{S_{\rm gen}}\,w_i^{\rm gen},
 \qquad \mathcal L=10^4\,\mathrm{pb}^{-1}.
\]
$\mathcal L$ 是目标积分亮度，即前述 $10\,\mathrm{fb}^{-1}$；$\sigma_{\rm eff}$ 是样本的有效截面；
$S_{\rm gen}$ 是官方参数表提供的生成器权重和，本项目直接采用表中数值。
官方固定版本参数表\cite{SampleInfo}给出的数值与项目配置一致：
\begin{center}
\small
\begin{tabular}{lrr}
\toprule
样本 & $\sigma_{\rm eff}$（$\mathrm{pb}$） & $S_{\rm gen}$ \\
\midrule
信号 & 0.0060239 & 27881776.6536 \\
背景 & 1.2578 & 7538705.8077 \\
\bottomrule
\end{tabular}
\end{center}

\textbf{第二步：补偿只使用最终评估样本的比例。}
第 2.2 节将样本分为训练、验证和最终评估三部分，最终评估的分配比例为 $20\%=0.2$。
因此计算产额时，将物理权重除以 0.2：
\[
 u_i=\frac{w_i^{\rm phys}}{0.2}.
\]
\textbf{第三步：对事件权重求和，得到筛选前后的产额。}
令 $E$ 为某一过程通过物理筛选的全部最终评估事件，$s_i$ 为网络分数，
$t$ 为用验证集背景确定的分数阈值，具体方法见第~\ref{sec:threshold}~节。筛选前后的产额分别为
\[
 Y_{\rm before}=\sum_{i\in E}u_i,\qquad
 Y_{\rm high}=\sum_{i\in E,\,s_i\geq t}u_i.
\]
$Y_{\rm before}$ 对全部评估事件求和，是网络筛选前的产额；$Y_{\rm high}$ 只对达到分数阈值的事件求和，是筛选后的产额。
权重可以为负，求和时保留正负号。由于每个事件按权重贡献，很多条模拟事件也可能只对应很少的预期事件。

\textbf{本项目采用的归一化设置。}
项目不再另乘 $k$ 因子或过滤效率（配置均为 1）；有效截面中分支比与过滤效率的具体分解仍待生产说明核实，不能重复乘入。
当前也未乘堆积、轻子和触发尺度因子，与官方示例\cite{HZZAnalysis}有别，不能据此复现其产额。

\FloatBarrier

\section{特征组与神经网络方法}

\subsection{输入信息与学习目标}

先将网络筛选前的每个事件整理成一组数值，再把这些数值送入网络。
表~\ref{tab:feature-overview} 列出四组输入：$p_T$ 表示垂直于束流方向的动量，赝快度 $\eta$ 描述相对束流的方向，
$\Delta R$ 表示两个轻子在赝快度与方位角平面内的角分离。四组均从公开轻子的四动量取得或计算，
具体变量名、重建公式及角度约定见附录~\ref{app:fields}。
\begin{table}[H]
\caption{特征组与输入维数。ABCD、BCD、ACD、ABD、ABC 的输入维数依次为 19、11、15、17、14。}
\label{tab:feature-overview}
\centering\small
\begin{tabularx}{\linewidth}{@{}cp{.36\linewidth}Yc@{}}
\toprule
组 & 变量 & 物理含义 & 维数 \\
\midrule
A & 四个轻子的 $p_T$、$\eta$ & 单轻子的动量尺度与方向 & 8 \\
B & $m_{Z1}$、$m_{Z2}$、两对内的 $\Delta R$ & 轻子对质量与内部角分离 & 4 \\
C & $p_{T,4\ell}$、$|\Delta\phi_{ZZ}|$ & 四轻子整体横向运动与两对间方位角关系 & 2 \\
D & 三个角余弦及两个平面角 & 产生与衰变的空间取向 & 5 \\
\bottomrule
\end{tabularx}
\end{table}
每一种输入方案都单独训练。
所有分类器均不显式输入 $m_{4\ell}$，但其他变量也由这些轻子的四动量构造，仍可能与质量相关。

训练时已知每条模拟事件来自信号还是背景，分别赋予标签 1 和 0。
训练的目标是让网络给信号事件较高的分数、给背景事件较低的分数，以便按分数筛选事件。
输入网络前，先对每个变量作标准化。令 $x_{ij}$ 为事件 $i$ 的第 $j$ 个输入变量，
$\mu_j$、$\sigma_j$ 为该变量在训练集中的均值和标准差，则送入网络的数值为
\[
 \widetilde x_{ij}=\frac{x_{ij}-\mu_j}{\sigma_j}.
\]
这样可将单位和数值范围不同的变量换算到相近的数值尺度，便于训练。
验证和评估样本使用训练集得到的同一组 $\mu_j$、$\sigma_j$，不重新计算。
若训练集中某个变量的标准差为零，计算时将该分母设为 1。
下面先说明一个分数怎样算出，再说明网络参数怎样通过训练确定。

\subsection{三层隐藏层的全连接模型}\label{sec:network}

本研究使用 PyTorch 实现的\textbf{多层感知机}（MLP）。
每个事件对应一行运动学数值，这一行就是网络的输入向量。
“全连接”表示某层的每个节点都接收前一层所有节点的输出；
“前馈”表示信息从输入层依次传到输出层。

网络从输入层经过三个隐藏层，再到输出层。隐藏层是输入与输出之间的中间计算层，各层单元数依次为：
\[
 \underbrace{d_{\rm in}}_{\text{输入}}
 \;\longrightarrow\;
 \underbrace{64\;\longrightarrow\;64\;\longrightarrow\;32}_{\text{三个隐藏层}}
 \;\longrightarrow\;
 \underbrace{1}_{\text{输出}}.
\]
输入维数 $d_{\rm in}$ 就是一条事件送入网络的数值个数，由表~\ref{tab:feature-overview} 中的方案决定。
隐藏层中的 64 或 32 表示计算节点数，不是一次输入多少事件，也不是新增了多少物理变量。
每个中间节点都可组合前一层传来的数值，不预先指定某个节点只负责某一物理量或特征组。

\begin{table}[H]
\caption{实际网络的计算顺序。Linear 表示全连接线性变换；各方案只改变输入维数，隐藏层结构相同。}
\label{tab:network}
\centering\small
\begin{tabularx}{\linewidth}{@{}p{.15\linewidth}p{.18\linewidth}Y@{}}
\toprule
部分 & 维数变化 & 依次执行的操作 \\
\midrule
输入处理 & $d_{\rm in}\to d_{\rm in}$ & 使用训练集的均值和标准差，对每个输入变量作标准化。 \\
隐藏层 1 & $d_{\rm in}\to64$ & Linear、LayerNorm、SiLU、Dropout（0.1）。 \\
隐藏层 2 & $64\to64$ & Linear、LayerNorm、SiLU、Dropout（0.1）。 \\
隐藏层 3 & $64\to32$ & Linear、LayerNorm、SiLU。 \\
输出 & $32\to1$ & Linear 得到实数 $z$，计算分数时再施加 sigmoid。 \\
\bottomrule
\end{tabularx}
\end{table}

对一个节点来说，先将前一层各数值乘以对应权重并相加，再加上偏置：
\[
 a_j=\sum_k W_{jk}h_k+b_j.
\]
$h_k$ 是上一层传入的数值，$W_{jk}$ 决定该数值对节点 $j$ 的贡献，$b_j$ 是额外的偏置。
第一层的 $h_k$ 是标准化后的物理输入，后续层则接收上一隐藏层的输出。
训练会调整这些权重、偏置以及归一化层的可学习参数，使网络给信号和背景的分数更符合各自标签。

\Needspace{6\baselineskip}
表~\ref{tab:network} 中的另外三种操作可按其作用理解：
\begin{itemize}
\item \textbf{LayerNorm（层归一化）：} 对同一个事件在该隐藏层的一组数值作归一化，再作可学习的缩放和平移。
它与输入标准化不同：输入标准化的统计量来自训练集，LayerNorm 则在每个事件内部计算，
不需要用最终评估样本重新拟合一套均值和标准差。
\item \textbf{SiLU（非线性激活）：} 将输入 $a$ 变为 $a/(1+e^{-a})$。
这使网络可以学习非线性的变量组合；如果所有层只做线性变换，多层仍只能合并为一个线性变换。
\item \textbf{Dropout（训练时随机屏蔽）：} 前两个隐藏层在训练时以 0.1 的概率将某个输出置零，
其余输出按 $1/0.9$ 缩放，用来减少对少数中间节点的依赖。
验证与最终评估时关闭随机屏蔽，使用完整网络；并非删除 10\% 的事件或物理特征组。
\end{itemize}

输出层产生一个可正可负的数 $z$，随后用
\[
 s=\frac{1}{1+e^{-z}}
\]
得到 $0$--$1$ 范围的分数 $s$。高分表示网络将该事件排在更接近信号的一侧，
分数本身不是产额、质量或信号强度；产额还需在施加阈值后，对留下事件的物理权重求和。
训练使用了类别内重标度的权重，因此这个分数不能直接当作真实信号比例或已校准的信号概率。

\subsection{网络训练和方案对比}

训练的目标是调整网络参数，让信号事件得到较高的分数、背景事件得到较低的分数。
下面按一批事件的处理过程说明。

\textbf{第一步：取出一批事件，计算分数。}
每轮先随机打乱训练事件的顺序，再按最多 1024 个事件组成一批。
将这一批事件的标准化输入送入当前网络，按第~\ref{sec:network}~节的计算得到每个事件的分数 $s_i$。
此时分数由当前网络参数决定，每个事件的信号或背景标签已经由模拟样本给出。

\textbf{第二步：用分数和标签计算分类损失。}
令 $y_i=1$ 表示信号，$y_i=0$ 表示背景，将已知标签与当前分数代入二元交叉熵公式：
\[
 \ell_i=-\left[y_i\ln s_i+(1-y_i)\ln(1-s_i)\right].
\]
对于信号，公式变为 $\ell_i=-\ln s_i$，分数越低，损失越大；
对于背景，公式变为 $\ell_i=-\ln(1-s_i)$，分数越高，损失越大。
因此，$\ell_i$ 用来衡量当前分数是否符合标签。训练通过减小损失，促使信号分数升高、背景分数降低。

\textbf{第三步：加权汇总这一批的损失。}
每个事件对总损失的贡献由优化权重 $\alpha_i$ 决定。训练开始前，先用全部训练事件，
分别计算信号和背景的平均绝对物理权重，再得到各事件的优化权重：
\[
 \overline w_c=\frac{1}{N_c}\sum_{i\in\mathcal T_c}|w_i^{\rm phys}|,
 \qquad
 \alpha_i=\frac{|w_i^{\rm phys}|}{\overline w_{y_i}}.
\]
$c=0,1$ 分别表示背景和信号，$\mathcal T_c$ 是该类的全部训练事件，$N_c$ 是其中的事件数。
使用时，将事件的绝对物理权重除以其所属类别的平均值。
取绝对值可使损失权重非负；除以类别平均值，可使两类的平均单事件优化权重均为 1，
避免两类物理权重的整体尺度不同影响训练，同时保留同一类别内事件之间的相对权重。
两类的事件数可以不同，因此它们的总优化权重不一定相等。

令 $\mathcal B$ 为当前批次，$m$ 为该批的事件数。把第二步得到的损失乘以优化权重，再按事件数取平均：
\[
 L_{\rm batch}=\frac{1}{m}\sum_{i\in\mathcal B}\alpha_i\ell_i.
\]
这样得到一个代表整批分类情况的损失 $L_{\rm batch}$，用于下一步更新网络参数。
其中 $\alpha_i$ 服务于训练；计算最终产额时，仍使用第~\ref{sec:weights}~节的带符号物理权重。

\textbf{第四步：根据损失更新网络参数。}
反向传播计算 $L_{\rm batch}$ 对各网络参数的导数，说明参数变化会怎样影响损失。
AdamW 优化器利用这些导数更新网络的连接权重、偏置等可学习参数。
本研究的学习率为 $10^{-3}$，用于控制更新步长的尺度；权重衰减为 $10^{-4}$，在更新时对参数作收缩，抑制参数过度增大。

\textbf{第五步：继续处理下一批，并重复多轮。}
更新参数后，取下一批事件，用更新后的网络重新计算分数和损失，再更新参数。
遍历全部训练事件一次称为一轮（epoch）；下一轮重新打乱事件顺序，重复上述过程。
每轮结束后用验证集检查模型表现，决定是否保存当前参数以及是否继续训练，具体规则见下一节。

\textbf{不同输入方案怎样比较。}
五种输入方案都按上述流程单独训练，使用相同的隐藏层结构、优化器、训练轮数上限、样本分工和检查点选择规则。
它们改变的是输入包含哪些特征；输入维数的变化也会改变第一层的参数数量。
每种方案分别使用种子 42、43、44，从随机初始化开始训练，共得到十五个模型。
种子控制参数初始化、事件打乱和 Dropout 的随机性，所有方案和种子使用的物理事件划分保持相同。
后文先比较同一种子下不同方案的结果，再汇总三个种子的差异。

\Needspace{15\baselineskip}
\subsection{用验证 AUC 选择检查点}

训练每完成一轮，都会得到一组网络参数。本节说明怎样从这些参数中选择最终使用的一组，即检查点。

\textbf{第一步：在验证集上检查分类表现。}
每轮训练后，用当前网络计算验证事件的分数，再用绝对物理权重计算验证 AUC，作为选择检查点的依据。
这里用验证集检查模型对未参与参数更新的事件表现如何。

\textbf{第二步：保存有改进的检查点。}
只有当前验证 AUC 比此前保存的最佳值提高超过 $10^{-4}$，才保存当前网络参数，并更新最佳值。
因此，最终选择依据是验证 AUC，而不是训练结束时的最后一组参数。

\textbf{第三步：停止训练，恢复已保存的最佳参数。}
连续 20 轮没有上述改进，就提前停止；训练最多进行 200 轮。停止后恢复保存的最佳检查点。
例如种子 42 的 ABCD 在第 87 轮结束训练，最终使用第 67 轮的参数；BCD 在第 86 轮结束，使用第 66 轮。
选定检查点后，再按下一节的方法确定分数阈值；这两步承担不同任务。
实现与复现说明见附录~\ref{app:network}。

\Needspace{13\baselineskip}
\subsection{用验证背景确定分数阈值}
\label{sec:threshold}

选定网络后，还需确定多高的分数才算通过筛选。这个分界分数就是阈值 $t$。
本研究用验证集背景的加权中位数确定阈值，使选择规则在最终评估前固定下来。
每个模型只使用自己的\emph{验证集背景分数}，依次进行：
\begin{enumerate}[leftmargin=1.7em,itemsep=2pt,topsep=3pt]
\item 将验证背景按网络分数从小到大排序。
\item 累加每个事件的绝对物理权重，找到累计权重第一次达到总权重一半的位置。
\item 将该位置的实际分数记为阈值 $t$；不在相邻分数之间插值。
\end{enumerate}
\textbf{怎样使用阈值。} 对每个事件比较分数与阈值：$s\geq t$ 划入高分类别，$s<t$ 划入低分类别。
分数恰好等于阈值的事件全部进入高分类别，不拆分同分事件。阈值确定后，不再根据最终评估结果调整。

\textbf{怎样比较不同模型。} 种子 42 的阈值分别为 ABCD 的 $0.746724$、BCD 的 $0.860891$。
两个网络的分数刻度不同，阈值更大不表示筛选更好。
应在相同评估事件上应用各自阈值，再比较留下的信号和背景。

\Needspace{15\baselineskip}
\subsection{冻结后应用于最终评估样本}

最终评估用于检查已经确定的筛选方法。开始前，先固定输入标准化参数、网络检查点和分数阈值，
这一步称为“冻结”。随后按以下顺序处理评估事件：
\begin{enumerate}[leftmargin=1.7em,itemsep=2pt,topsep=3pt]
\item 用训练集保存的标准化参数处理输入，再用选定的网络计算分数。
\item 将分数与该模型的固定阈值比较，分为高分和低分两类。
\item 分别统计事件数、求带符号评估权重之和，并绘制质量谱。
\end{enumerate}
本报告的“网络筛选后”指高分类别。事件质量本身保持原值，质量谱的变化来自留下了哪些事件。

\textbf{为什么背景不一定正好保留一半。}
阈值虽然按验证背景的加权中位数确定，但最终评估的背景产额不一定保留 50\%。
一是事件权重和同分事件不能任意拆分，验证集上的高分绝对权重占比就可能偏离一半；
二是评估集与验证集包含不同事件；三是最终产额使用带符号权重，而阈值使用绝对权重。
附录~\ref{app:threshold} 给出一个明确标为教学示例的计算。

基础指标使用整个 $[105,140)$ GeV 范围内的高分产额。
同时保留低分产额，并检查高、低两类产额之和是否等于网络筛选前的总产额，核对事件是否完整分配。

\section{比较指标与计算示例}\label{sec:basic}

\subsection{事件数与加权产额}
先比较种子 42 的两个模型：ABCD 使用全部输入，BCD 去掉单轻子变量 A。
表~\ref{tab:groups} 先列出它们各留下多少物理事件组，同一物理事件只计一次。
例如，ABCD 从 9,049 个信号事件组中保留了 8,761 个。

这些数字说明实际使用了多少模拟事件。要得到预期产额，还需按第~\ref{sec:weights}~节对带符号权重求和。
\begin{table}[H]
\caption{本次运行最终评估样本中的物理事件组数。后两列为种子 42 的高分子集。}
\label{tab:groups}
\centering
\begin{tabular}{lrrr}
\toprule
过程 & 网络筛选前 & ABCD 高分 & BCD 高分 \\
\midrule
信号 & 9,049 & 8,761 & 8,563 \\
背景 & 698 & 386 & 378 \\
\bottomrule
\end{tabular}
\end{table}

\subsection{从产额到保留率和信背比}
本节用三个指标回答：信号保留了多少、背景保留了多少，以及筛选后信号相对于背景有多少。

\textbf{第一步：计算产额和三项指标。}
令 $E_S$、$E_B$ 分别为网络筛选前的最终评估信号、背景事件集合，$u_i$ 为第~\ref{sec:weights}~节的评估权重。
将各集合内全部事件的权重相加，得到网络筛选前的信号产额 $S_0$ 和背景产额 $B_0$：
\[
 S_0=\sum_{i\in E_S}u_i,\qquad
 B_0=\sum_{i\in E_B}u_i.
\]
令 $s_i$ 为事件分数，$t$ 为该模型已固定的阈值。只对通过阈值的事件求和，得到网络筛选后的产额：
\[
 S=\sum_{i\in E_S,\,s_i\geq t}u_i,\qquad
 B=\sum_{i\in E_B,\,s_i\geq t}u_i.
\]
这四项产额均对应同一 $10\,\mathrm{fb}^{-1}$。再将它们代入指标公式：
\[
 R_S=\frac{S}{S_0},\qquad R_B=\frac{B}{B_0},\qquad Q=\frac{S}{B}.
\]
$R_S$ 将筛选后的信号产额除以筛选前的信号产额，得到信号保留率；$R_B$ 对背景作同样计算。
背景去除比例为 $1-R_B$。$Q$ 则直接比较筛选后的信号与背景，例如 $Q=2$ 表示信号产额为背景的两倍。
保留率以百分数列出，产额和信背比直接列数值。

将 $S=R_SS_0$、$B=R_BB_0$ 代入信背比，可得
\[
 Q=\frac{S_0}{B_0}\frac{R_S}{R_B}.
\]
这个式子说明：网络筛选前的产额固定时，信背比由信号和背景的保留比例共同决定。
例如，背景减少得更多时，信背比可能提高，但同时也可能损失信号。因此需要一起报告三项指标，
不能只按 $Q$ 排名或合成一个总分。

各模型使用同样的阈值确定规则，但评估时的背景保留率不一定相同。
这里比较各模型按固定阈值留下的实际产额，没有再调阈值使它们保留相同比例的信号或背景。

\textbf{第二步：检查分母和带符号结果。}
若 $S_0$ 或 $B_0$ 非正，对应保留率不计算；若 $B\leq0$，信背比不计算，并说明原因。
产额始终使用带符号权重，不能先取绝对值。负产额或超出 $[0,1]$ 的保留率须保留原值并标记解释限制。

\subsection{ABCD 与 BCD 的计算示例}
\begin{table}[H]
\caption{本次运行在 $10\,\mathrm{fb}^{-1}$ 下的预期带符号产额及保留率。
误差为固定网络和阈值下按物理事件组计算的 MC 标准误差；保留率误差包括筛选前后共享事件的协方差。}
\label{tab:yields}
\centering
\begin{tabular}{llrr}
\toprule
过程 & 选择 & 预期产额 & 保留率（\%） \\
\midrule
信号 & 网络筛选前（$S_0$） & $2.7710\pm0.0293$ & --- \\
 & ABCD 高分 & $2.6824\pm0.0288$ & $96.80\pm0.19$ \\
 & BCD 高分 & $2.6215\pm0.0285$ & $94.60\pm0.24$ \\
背景 & 网络筛选前（$B_0$） & $2.3553\pm0.1683$ & --- \\
 & ABCD 高分 & $1.2027\pm0.1202$ & $51.06\pm3.57$ \\
 & BCD 高分 & $1.1012\pm0.1156$ & $46.75\pm3.57$ \\
\bottomrule
\end{tabular}
\end{table}

将表~\ref{tab:yields} 的产额代入上一节的公式，得到表~\ref{tab:worked-metrics}。
最后一列按“BCD 减 ABCD”计算，比较去掉 A 后各项指标的变化。
\begin{table}[H]
\caption{种子 42 的指标计算与删去 A 后的变化。这里使用显示精度演示计算；完整结果由未舍入汇总计算。
保留率差以百分点表示，信背比差为直接差值。}
\label{tab:worked-metrics}
\centering\small
\renewcommand{\arraystretch}{2.6}
\begin{tabular}{lccr}
\toprule
指标 & ABCD & BCD & BCD $-$ ABCD \\
\midrule
信号保留率 $R_S$ & $\dfrac{2.6824}{2.7710}\simeq96.80\%$ & $\dfrac{2.6215}{2.7710}\simeq94.60\%$ & $-2.20$ 个百分点 \\
背景保留率 $R_B$ & $\dfrac{1.2027}{2.3553}\simeq51.06\%$ & $\dfrac{1.1012}{2.3553}\simeq46.75\%$ & $-4.31$ 个百分点 \\
信背比 $Q$ & $\dfrac{2.6824}{1.2027}\simeq2.23$ & $\dfrac{2.6215}{1.1012}\simeq2.38$ & 约 $+0.15$ \\
\bottomrule
\end{tabular}
\end{table}
在种子 42 的示例中，BCD 留下更少信号和背景，信背比提高。
这体现了信号损失与背景减少之间的取舍，不能由这一次比较判断 BCD 普遍更优。

\Needspace{12\baselineskip}
\textbf{理解表中的误差。}
表~\ref{tab:yields} 中 $\pm$ 后的数值表示有限 MC 样本带来的标准误差。
对当前过程（信号或背景）和当前选择，令 $\mathcal G$ 为参与计算的全部物理事件组的集合，$g$ 为事件组编号，
$\mathcal I_g$ 为第 $g$ 组中参与当前产额计算的模拟记录集合，$u_i$ 为记录 $i$ 的带符号评估权重。
先计算每组的贡献 $U_g$，再按现有事件组计数约定计算产额和标准误差：
\[
 U_g=\sum_{i\in\mathcal I_g}u_i,\qquad
 Y=\sum_{g\in\mathcal G}U_g,\qquad
 \sigma_Y=\sqrt{\sum_{g\in\mathcal G}U_g^2}.
\]
同一物理事件对应的记录归为一组，先合并组内贡献，再计算误差，避免将相关记录当作独立事件。
只有独立单位权重计数等特殊情形，误差才化为 $\sqrt{N}$。
这里固定了网络和阈值，误差不包含重复训练的波动。保留率的分子与分母还共享事件，
计算保留率误差时需计入协方差，具体方法见附录~\ref{app:errors}。


\Needspace{14\baselineskip}
\subsection{删组比较：先做差，再汇总}
本节比较去掉一组输入后，指标相对 ABCD 怎样变化。去掉 A、B、C、D 分别得到 BCD、ACD、ABD、ABC。

\textbf{第一步：在同一个种子内作差。}
令 $G$ 为去掉的特征组，$r$ 为种子。将删组方案的指标减去同种子 ABCD 的指标：
\[
 \Delta R_{S,G,r}=R_S(\mathrm{ABCD}\setminus G,r)-R_S(\mathrm{ABCD},r),
\]
\[
 \Delta R_{B,G,r}=R_B(\mathrm{ABCD}\setminus G,r)-R_B(\mathrm{ABCD},r),\qquad
 \Delta Q_{G,r}=Q(\mathrm{ABCD}\setminus G,r)-Q(\mathrm{ABCD},r).
\]
使用公式时，始终按“删组方案减 ABCD”的顺序。两项保留率差乘以 100 后以\textbf{百分点}报告；
$\Delta Q$ 则直接报告信背比差。

\textbf{第二步：分别解释三个差值。}
负的 $\Delta R_S$ 表示信号保留得更少；负的 $\Delta R_B$ 表示背景保留得更少；正的 $\Delta Q$ 表示信背比提高。
因此，不能只看正负号判断整体好坏，需要同时说明信号损失和背景减少。

\textbf{第三步：汇总三个种子的差值。}
对同一删组和同一指标，将三个种子的差值排序，取中间值为中位数，并列出最小值和最大值。
这样先保留每个种子内的配对比较，再观察训练随机性带来的变化。全部三个种子都要纳入汇总。
先分别取两种方案的中位数、再相减，一般不等于上述方法，因此本文统一先作差再汇总。

三个差值同号，表示这三次训练的变化方向相同；有正有负，则表示方向随种子改变。
最小至最大范围仅描述这三个差值，不是置信区间，不能用来判断统计显著性。
此外，删组方案重新训练了网络，也重新确定了阈值。差值反映整套筛选方法的变化，不能解释为该组独立的物理贡献。


\FloatBarrier
\Needspace{12\baselineskip}
\section{五种输入方案的比较结果}

\subsection{逐种子的基础指标}
表~\ref{tab:basic-all} 列出五种输入方案、每种三个种子的全部十五个模型。
各模型使用相同的网络筛选前产额：$S_0=2.7710$、$B_0=2.3553$，且高、低两类的产额之和均与网络筛选前的产额一致。
因此，表中差别来自各模型筛选后留下的事件。

读表时，先选定一个种子，比较五种方案的信号保留率 $R_S$ 和背景保留率 $R_B$，再看信背比 $Q$。
随后检查换一个种子时，变化方向是否仍然相同。全部十五个模型均列出，避免只挑选表现最好的一次训练。
\begin{table}[H]
\caption{本机新运行的基础指标表。全部十五个模型的汇总均有效，数值由未舍入的带符号产额计算。
$R_S$、$R_B$ 的单位为 \%，$S$、$B$ 对应 $10\,\mathrm{fb}^{-1}$。}
\label{tab:basic-all}
\centering\small
\begin{tabular}{llrrrrr}
\toprule
输入 & 种子 & $S$ & $B$ & $R_S$ & $R_B$ & $Q=S/B$ \\
\midrule
ABCD & 42 & 2.6824 & 1.2027 & 96.80 & 51.06 & 2.230 \\
ABCD & 43 & 2.6809 & 1.1589 & 96.75 & 49.20 & 2.313 \\
ABCD & 44 & 2.6670 & 1.0834 & 96.25 & 46.00 & 2.462 \\
\addlinespace
BCD & 42 & 2.6215 & 1.1012 & 94.60 & 46.75 & 2.381 \\
BCD & 43 & 2.6156 & 1.1145 & 94.39 & 47.32 & 2.347 \\
BCD & 44 & 2.6264 & 1.1344 & 94.78 & 48.16 & 2.315 \\
\addlinespace
ACD & 42 & 2.6667 & 1.1536 & 96.24 & 48.98 & 2.312 \\
ACD & 43 & 2.6679 & 1.1232 & 96.28 & 47.69 & 2.375 \\
ACD & 44 & 2.7039 & 1.2813 & 97.58 & 54.40 & 2.110 \\
\addlinespace
ABD & 42 & 2.5731 & 1.2186 & 92.86 & 51.74 & 2.112 \\
ABD & 43 & 2.5408 & 1.2192 & 91.69 & 51.76 & 2.084 \\
ABD & 44 & 2.5805 & 1.2026 & 93.13 & 51.06 & 2.146 \\
\addlinespace
ABC & 42 & 2.6975 & 1.1241 & 97.35 & 47.73 & 2.400 \\
ABC & 43 & 2.6972 & 1.0311 & 97.33 & 43.78 & 2.616 \\
ABC & 44 & 2.6941 & 1.1386 & 97.22 & 48.34 & 2.366 \\
\bottomrule
\end{tabular}
\end{table}

\subsection{删组差值与种子变化}
表~\ref{tab:basic-deltas} 按第~\ref{sec:basic}~节的方法，将上一表转换为“删组方案减 ABCD”的差值。
先读每个种子的数值，再看中位数和最小至最大范围。图~\ref{fig:basic-deltas} 展示同一组差值：
点位于零线的哪一侧表示差值的正负，可直接检查三个种子的方向是否一致。
\begin{table}[H]
\caption{本机新运行的删组差值汇总。保留率差以百分点表示，信背比差无单位；每个汇总必须包含全部三个种子。}
\label{tab:basic-deltas}
\centering\small
\begin{tabular}{llccccc}
\toprule
去掉组 & 指标 & 种子 42 & 种子 43 & 种子 44 & 中位数 & 最小至最大 \\
\midrule
A & $\Delta R_S$ & $-2.20$ & $-2.35$ & $-1.47$ & $-2.20$ & $[-2.35,-1.47]$ \\
A & $\Delta R_B$ & $-4.31$ & $-1.88$ & $+2.17$ & $-1.88$ & $[-4.31,+2.17]$ \\
A & $\Delta Q$ & $+0.150$ & $+0.033$ & $-0.147$ & $+0.033$ & $[-0.147,+0.150]$ \\
\addlinespace
B & $\Delta R_S$ & $-0.57$ & $-0.47$ & $+1.33$ & $-0.47$ & $[-0.57,+1.33]$ \\
B & $\Delta R_B$ & $-2.09$ & $-1.51$ & $+8.40$ & $-1.51$ & $[-2.09,+8.40]$ \\
B & $\Delta Q$ & $+0.081$ & $+0.062$ & $-0.351$ & $+0.062$ & $[-0.351,+0.081]$ \\
\addlinespace
C & $\Delta R_S$ & $-3.94$ & $-5.05$ & $-3.12$ & $-3.94$ & $[-5.05,-3.12]$ \\
C & $\Delta R_B$ & $+0.68$ & $+2.56$ & $+5.06$ & $+2.56$ & $[+0.68,+5.06]$ \\
C & $\Delta Q$ & $-0.119$ & $-0.229$ & $-0.316$ & $-0.229$ & $[-0.316,-0.119]$ \\
\addlinespace
D & $\Delta R_S$ & $+0.54$ & $+0.59$ & $+0.98$ & $+0.59$ & $[+0.54,+0.98]$ \\
D & $\Delta R_B$ & $-3.34$ & $-5.42$ & $+2.34$ & $-3.34$ & $[-5.42,+2.34]$ \\
D & $\Delta Q$ & $+0.169$ & $+0.302$ & $-0.096$ & $+0.169$ & $[-0.096,+0.302]$ \\
\bottomrule
\end{tabular}
\end{table}

\begin{figure}[H]
\centering
\includegraphics[width=\linewidth]{figures/basic-20261108/deletion-differences.pdf}
\caption{本机新运行的逐种子删组差值。横轴为去掉的输入组；三面板依次为信号保留率差（百分点）、
背景保留率差（百分点）和信背比差。圆点表示种子 42--44，空心菱形表示中位数，
细线连接最小值和最大值，虚线为零线。细线不是置信区间。}
\label{fig:basic-deltas}
\end{figure}

\subsection{结果分析}
\textbf{删去 C。} 三个种子的信号保留率均下降（$-5.05$ 至 $-3.12$ 个百分点），
背景保留率均上升（$+0.68$ 至 $+5.06$ 个百分点），信背比均下降（$-0.316$ 至 $-0.119$）。
这三次训练中，保留 C 组的 ABCD 都比删去 C 的 ABD 留下更多信号、更少背景，并得到更高信背比。
这一结果适用于各自已冻结的网络和阈值。C 组包括四轻子整体横动量及两个轻子对间的方位角差，
本比较尚不能区分这两个变量各自起了多大作用。

\textbf{删去 A。} 信号保留率均下降（$-2.35$ 至 $-1.47$ 个百分点），
但背景保留率和信背比的变化方向不一致。种子 42、43 留下更少信号和更少背景，
信背比提高；种子 44 则信号减少、背景增加，信背比降低。因此，删去 A 的结果随种子改变。

\textbf{删去 B。} 三项指标的差值均存在种子间符号变化。
种子 42、43 的信号和背景保留率均略降，信背比提高；种子 44 则保留更多信号和背景，信背比下降。
因此，读这一行的中位数时，还要同时查看三个种子的正负号，避免把不同方向的变化概括成同一种效果。

\textbf{删去 D。} 信号保留率均上升（$+0.54$ 至 $+0.98$ 个百分点），
但背景保留率和信背比没有一致改善。种子 42、43 的背景减少、信背比提高，种子 44 的方向相反。

\Needspace{4\baselineskip}
综合三项指标，只有删去 C 在三个种子中都表现为信号减少、背景增加、信背比降低。
其余删组至少有一项指标的方向随种子改变，需逐项解释。
这些结果限于当前样本、训练设置和阈值规则，尚不能确定普遍最优的输入组合、独立特征贡献或差异的统计显著性。
\FloatBarrier

\Needspace{14\baselineskip}
\section{质量谱变化}
保留率说明总产额留下了多少，质量谱则说明这些产额分布在哪些质量区域。
本节比较同一评估样本上、种子 42 的 ABCD 和 BCD：先看各区域的绝对产额，再比较它们占总产额的比例。
这样可区分“留下的总量变了”和“留下事件的质量分布变了”。
这里仅检查这两个模型，不能据此解释全部删组结果或确定 C 组的作用机制。
展示时使用七个 5 GeV 质量箱；前面各项基础指标仍按整个质量窗口计算。

\subsection{绝对产额谱：信号和背景位于哪里}
\textbf{怎样计算。} 按事件的四轻子不变质量分箱，再分别计算信号和背景在每个箱中的产额。
令 $b$ 表示一个质量箱，用第~\ref{sec:weights}~节的评估权重计算
\[
 Y_b=\sum_{i\in b}u_i.
\]
先对全部评估事件计算一条筛选前曲线，再只对通过网络阈值的事件计算筛选后曲线。
各箱的产额相加，就是对应选择在整个质量窗口中的产额。

\textbf{怎样读图。} 图~\ref{fig:yields} 左列为信号，右列为背景。
比较同一面板中的筛选前后曲线，可看到各质量区域留下多少预期产额。
两条曲线使用相同积分亮度和质量箱边界；信号与背景面板的纵轴刻度不同。

\begin{figure}[H]
\centering
\includegraphics[width=.92\linewidth]{figures/basic-20261108/mass-yields.pdf}
\caption{本次运行最终评估的带符号质量谱。上排为 ABCD，下排为 BCD，均使用种子 42。
蓝色圆点表示网络筛选前的产额，橙色圆点表示高分事件的产额；左列为信号，右列为背景。每点代表一个 5 GeV 箱，
连线仅作视觉引导，不代表连续质量模型。纵轴为 $10\,\mathrm{fb}^{-1}$ 下的预期产额，信号与背景刻度不同。
误差棒表示按物理事件组计算的 MC 标准误差。阴影标记统计支持不足的区域；负产额保留。}
\label{fig:yields}
\end{figure}



\textbf{信号结果。} 模拟信号集中在 125 GeV 附近，ABCD 和 BCD 的筛选后曲线都基本保留这一主要区域。
网络只决定事件是否留下，不改变事件质量。这里的粗分箱结果不能用来判断峰位是否移动、
探测器分辨率是否改善或希格斯物理宽度是否变化。

\textbf{背景结果。} 筛选前的背景覆盖整个质量窗口，分布本来就不均匀。
筛选后，各质量区域减少的程度不同，高质量尾部减少尤其明显。
为进一步比较剩余背景的分布，下一节将每条曲线除以自己的总产额。
\FloatBarrier

\Needspace{12\baselineskip}
\subsection{归一化谱形：去除总量变化后的比较}
绝对产额减少时，曲线整体会降低，单看高度不容易比较形状。
因此，图~\ref{fig:shapes} 将每条曲线分别归一化，只比较各质量区域占该曲线总产额的份额。

\textbf{怎样计算。} 对某一过程和某一种选择，令各箱产额为 $Y_b$，总产额为 $Y_{\rm total}=\sum_bY_b$。
在总产额为正时，计算
\[
 f_b=\frac{Y_b}{Y_{\rm total}}.
\]
$f_b$ 表示这一箱占该曲线总产额的比例。筛选前曲线用筛选前总产额作分母，筛选后曲线用筛选后总产额作分母。
使用带符号产额时，个别箱的份额仍可能为负，归一化曲线不一定是概率密度。
表~\ref{tab:fractions} 再将指定质量区域内各箱的份额相加，以百分数列出，便于具体比较。

\begin{figure}[H]
\centering
\includegraphics[width=.82\linewidth]{figures/basic-20261108/mass-shapes.pdf}
\caption{本次运行：与图~\ref{fig:yields} 使用相同事件和阈值、分别归一化后的筛选前后质量谱形。
每个箱除以其曲线自身为正的带符号总产额。误差传播包括共享归一化；筛选前后也共享事件。
带符号归一化谱形不一定是概率密度。保留阴影箱和负值箱，不作平滑、裁剪或重新分箱。}
\label{fig:shapes}
\end{figure}

\begin{table}[H]
\caption{本次运行固定质量区域内的带符号产额占各自曲线总产额的比例（\%）。
数值为描述性点估计，未展示不确定度，不能据此判定筛选前后差异的统计显著性。}
\label{tab:fractions}
\centering
\small
\begin{tabular}{llrrr}
\toprule
过程 & 区域（GeV） & 筛选前 & ABCD 高分 & BCD 高分 \\
\midrule
信号 & $[120,130)$ & 87.9 & 88.6 & 88.5 \\
背景 & $[120,130)$ & 32.5 & 51.8 & 47.9 \\
背景 & $[135,140)$ & 22.0 & 6.5 & 9.0 \\
\bottomrule
\end{tabular}
\end{table}

\textbf{怎样解释结果。} 信号仍集中于主要质量区域。背景在该区域的相对占比上升，高质量尾部的占比下降。
读这些比例时，要同时记住筛选后总背景已经减少：某一区域的占比上升，不代表该区域的绝对产额增加。

即使没有显式输入 $m_{4\ell}$，筛选仍改变了不同质量区域的相对份额，说明这一选择与质量有关。
这不能确定是哪一个输入变量造成了变化，也不能据此判断背景存在物理共振。
\FloatBarrier

\subsection{统计支持与解释限制}
\textbf{先检查阴影区域。} 阴影提醒读者：该箱事件组不足、产额估计不稳定，或正负权重抵消较强，
判定规则见附录~\ref{app:errors}。ABCD 高分背景在 $[105,110)$ GeV 箱中的产额为负；
若干背景箱及信号高质量尾部也缺少足够统计支持。这些结果按原值展示，负产额未裁剪为零。

\textbf{再理解误差棒。} 误差棒描述网络和阈值固定时，有限 MC 样本造成的标准误差。
同一事件可能同时参与筛选前后曲线；归一化时，各箱又共用同一总产额。
因此，图中的点及两条曲线存在相关性，不能把筛选前后当作两组独立测量。

\section{结论与局限}\label{sec:conclusion}

\textbf{本报告回答的问题。} 去掉一组输入、重新训练网络并按既定规则确定阈值后，信号和背景各留下多少？
为回答这个问题，本文比较五种输入方案、每种三个种子的信号保留率、背景保留率和信背比。

\textbf{主要结果。} 删去 C 组时，三个种子都出现信号保留率下降、背景保留率上升、信背比下降。
其余删组没有在三个种子上同时改善全部指标，应分别查看信号损失、背景减少及变化方向。
ABCD 和 BCD 的质量谱还显示：总背景减少的同时，剩余背景在各质量区域的相对份额也会变化。
因此，评价筛选效果时，应一起检查保留率、信背比和质量分布。

\textbf{这些结果能说明什么。} 结论适用于本次两份受控 MC 样本、既定网络和冻结阈值。
特征组之间可能相关，删组方案又重新训练了网络并重新确定阈值，所以差值描述整套筛选方法的变化，
不能分解为某组独立的物理贡献。

\textbf{不确定度包含什么。} 三个种子的最小至最大范围描述这几次训练的结果变化；
有限 MC 标准误差描述网络和阈值固定时的样本波动。这两种展示均未涵盖完整训练与阈值估计的不确定性。
本文尚未建立差异的统计显著性，也未证明测量精度提高。

\textbf{研究范围。} 本研究仍属探索性、非独立分析。重新划分或重新运行已使用的源样本，不会得到独立验证样本。
样本生产配置、发布前筛选、部分选择条件的外部依据和物理权重适用性仍待核实。
本项目用于教育与技术研究，不构成 ATLAS/CMS 官方结果、希格斯发现或物理测量。

\clearpage
\appendix

\section{字段、重建与输入约定}\label{app:fields}

\subsection{模拟与发布范围}
\textbf{公开文件提供什么。} 一般模拟过程依次包括：生成硬散射、模拟簇射与强子化、模拟探测器响应、
重建电子和缪子等对象，再进行预筛选并发布 ROOT 文件。
本项目从公开文件中的已重建轻子开始，读取其字段并构造候选；下文的“重建”指组合这些轻子的四动量，
不使用生成器真值四动量。

\textbf{哪些信息仍未确认。} 两份样本的生成器版本、簇射和探测器模拟配置仍需核对生产记录，不能直接采用参考论文中的设置。
CERN 记录确认发布前进行了宽松的对象与事件选择，并保留至少四个电子或缪子；具体动量、接受度及品质阈值尚未确认。
因此，本项目是在已经预筛选的公开样本上继续筛选。表~\ref{tab:population} 的计数不能用于计算相对全部生成事件的接受效率。
\subsection{读取字段}
\textbf{怎样读取事件。} 文件中的 \texttt{mini} 树按记录存放事件。一个 entry 对应一条记录，branch 对应记录中的一个字段。
下表列出本项目读取的全部 19 个分支；原文件可能有其他分支，这 19 个分支也不等于网络的 19 维输入。

前七个字段每条记录只有一个数值，称为标量；其余字段为每个轻子保存一个数值，合起来称为向量。
例如 \texttt{lep\_n=4} 时，各轻子向量有四个元素。在不同向量中，相同下标对应同一个轻子，读取时据此组合它的各项信息。

\begingroup
\small
\renewcommand{\arraystretch}{1.13}
\begin{tabularx}{\linewidth}{@{}p{.62\linewidth}Y@{}}
\toprule
读取的分支 & 含义与存储单位 \\
\midrule
\texttt{runNumber, eventNumber, channelNumber} & 事件／样本标识（整数） \\
\texttt{mcWeight; trigE, trigM; lep\_n} & 生成器权重；触发布尔值；轻子数 \\
\texttt{lep\_pt, lep\_eta, lep\_phi, lep\_E} & $p_T$、$\eta$、$\phi$、能量；动量和能量为 MeV，角度为弧度 \\
\texttt{lep\_charge, lep\_type} & 电荷和粒子类型（整数向量） \\
\texttt{lep\_isTightID, lep\_trigMatched} & 鉴别和触发匹配标志（布尔向量） \\
\texttt{lep\_ptcone30, lep\_etcone20} & 径迹／量能器隔离量（MeV） \\
\texttt{lep\_z0; lep\_tracksigd0pvunbiased} & $z_0$（mm）；$d_0/\sigma(d_0)$（无量纲） \\
\bottomrule
\end{tabularx}
\endgroup

运动学、隔离及碰撞参数字段为浮点向量，使用前将 MeV 换算为 GeV。
事件权重按第~\ref{sec:weights}~节换算为预期产额。
这些字段说明来自项目输入配置和此前的 ROOT 元数据核对，尚不能补足样本生产过程的缺失信息。
\subsection{四动量重建与输入分组}
\textbf{第一步：构造单轻子四动量。} 对通过品质筛选的轻子，将动量和能量换算为 GeV，再按 $p_T$ 从高到低编号。
取 $c=1$，将公开的 $(p_T,\eta,\phi,E)$ 代入
\[
 p^\mu=(E,p_x,p_y,p_z),\qquad
 p_x=p_T\cos\phi,\quad p_y=p_T\sin\phi,\quad p_z=p_T\sinh\eta.
\]
这样得到每个轻子的能量和三个动量分量，后续可通过四动量相加计算组合系统的质量和动量。

\textbf{第二步：确定两个轻子对。} 按类型和电荷寻找同味异号（SFOS）对，即电子与正电子配对，或两个电荷相反的缪子配对。
本项目只研究 $2e2\mu$ 道，因此电子对和缪子对给出唯一的两对。
将质量更接近 $91.1876\gev$ 的一对记为 $Z_1$，另一对记为 $Z_2$。
文献中的 $m_{12},m_{34}$ 对应代码的 \texttt{mZ1,mZ2}；$Z_1$ 由质量确定，不固定对应电子对，也不按 $p_T$ 编号。

\textbf{第三步：计算组合系统的质量和横动量。}
先将每对轻子的四动量相加，得到 $P_j=p_{j^-}+p_{j^+}$，再计算四轻子总四动量 $Q=P_1+P_2$。
由此得到
\[
 m_{Zj}=\sqrt{P_j^2},\qquad m_{4\ell}=\sqrt{Q^2},\qquad
 p_{T,4\ell}=\sqrt{Q_x^2+Q_y^2}.
\]
前两项分别给出轻子对质量和四轻子质量，第三项给出四轻子系统垂直于束流方向的总动量大小。
这些计算使用已有轻子对象，不重新模拟探测器。所得质量谱仍受辐射、探测器响应和筛选影响，
谱的宽度不能直接解释为希格斯本征宽度。

\textbf{第四步：整理网络输入。} A 组直接取排序后的 $p_T,\eta$；B、C 组由上述四动量计算。
$\Delta R$ 使用同一对内两轻子的 $\eta,\phi$，方位角差先折回 $[-\pi,\pi)$；
\texttt{deltaPhi\_ZZ} 则是两个轻子对方位角差的绝对值。
下表的 19 维输入由读取的字段整理、计算而来，并非文件中另外存放的 19 个现成字段。
\begin{table}[H]
\caption{输入变量分组及其物理含义。动量和质量使用 GeV，方位角使用弧度。}
\label{tab:features}
\small
\renewcommand{\arraystretch}{1.12}
\begin{tabularx}{\linewidth}{@{}p{.05\linewidth}p{.43\linewidth}Y@{}}
\toprule
组 & 代码变量 & 物理含义 \\
\midrule
A & \texttt{lep1\_pt}--\texttt{lep4\_pt}；\texttt{lep1\_eta}--\texttt{lep4\_eta} &
单轻子的横动量和赝快度，描述动量尺度及相对束流的方向。 \\
B & \texttt{mZ1}, \texttt{mZ2}；\texttt{deltaR\_Z1}, \texttt{deltaR\_Z2} &
两个轻子对的质量，以及各对内部的角分离 $\Delta R=\sqrt{(\Delta\eta)^2+(\Delta\phi)^2}$。 \\
C & \texttt{pt4l}, \texttt{deltaPhi\_ZZ} &
四轻子系统的横动量，以及两个轻子对之间的方位角差绝对值。 \\
D & \texttt{cos\_theta\_star}, \texttt{cos\_theta\_1}, \texttt{cos\_theta\_2}；
\texttt{phi\_decay\_planes}, \texttt{phi\_production\_plane} &
一个产生角余弦、两个衰变极角余弦，以及描述衰变平面与产生平面相对取向的两个角。 \\
\bottomrule
\end{tabularx}
\end{table}
\subsection{五个角度的约定}
D 组在指定参考系中比较产生和衰变方向。计算前先对轻子四动量作洛伦兹变换，再按以下约定取角度：
\begin{enumerate}[leftmargin=1.7em,itemsep=3pt]
\item \textbf{产生角。} 在四轻子静止系中，比较 $Z_1$ 方向与变换后的实验室正束流方向，得到夹角余弦 $\cos\theta^*$。
\item \textbf{两个衰变极角。} 分别进入各轻子对静止系，比较负轻子方向与另一轻子对的反方向，得到 $\cos\theta_1,\cos\theta_2$。
\item \textbf{两个平面角。} 在四轻子静止系中，用同一对负、正轻子动量的叉积定义衰变平面法向量，
用正束流方向与 $Z_1$ 方向定义产生平面法向量。绕 $Z_1$ 轴取有向角，分别比较两个衰变平面，以及产生平面与第一衰变平面。
\end{enumerate}
角度取值范围为 $[-\pi,\pi)$，余弦为 $[-1,1]$。各输入方案使用相同约定，以保证变量含义一致。
\FloatBarrier

\Needspace{30\baselineskip}
\section{完整选择条件}\label{app:cuts}

\subsection{共同选择的执行顺序}
下表按执行顺序列出物理筛选条件及目的。每个事件在网络筛选前都须通过全部条件。
随后按第 2.2 节划分样本；最终评估中的网络筛选再使用已固定的分数阈值。
\begin{table}[H]
\caption{共同筛选的执行顺序。动量与质量单位为 GeV，分支含义见附录~\ref{app:fields}。}
\label{tab:cuts}
\small
\renewcommand{\arraystretch}{1.1}
\begin{tabularx}{\linewidth}{@{}p{.23\linewidth}Y@{}}
\toprule
步骤 & 实际条件及目的 \\
\midrule
\textbf{事件触发} & \texttt{trigE} 或 \texttt{trigM} 为真，要求事件通过电子或缪子触发。 \\
\midrule
\textbf{轻子品质} & 依次执行以下条件，每步至少保留四个轻子，以筛出符合要求的提示轻子候选： \\
味道、鉴别 & $|\texttt{lep\_type}|\in\{11,13\}$；\texttt{lep\_isTightID=true}。 \\
径迹、量能器隔离 & $p_T>0$；依次要求 \texttt{lep\_ptcone30}/\texttt{lep\_pt}$<0.3$、
\texttt{lep\_etcone20}/\texttt{lep\_pt}$<0.3$，抑制周围活动较多的对象。 \\
横向、纵向碰撞参数 & $|d_0/\sigma(d_0)|<$ 5（电子）或 3（缪子）；
$|z_0/\cosh\eta|<0.5$ mm，限制相对主顶点的位移。 \\
\midrule
\textbf{四轻子候选} & 以下条件依次定义共同分析总体： \\
多重数、触发匹配 & 品质筛选后恰好四个轻子；至少一个 \texttt{lep\_trigMatched=true}。
原始 \texttt{lep\_n} 可以大于四。 \\
排序动量、接受度 & 按 $p_T$ 降序，依次 $\geq20,15,10,7$；
电子 $|\eta|<2.47$、缪子 $|\eta|<2.7$，限定运动学及探测器接受范围。 \\
电荷、配对 & 总电荷为零；存在两个互不重叠的同味异号（SFOS）对，构造 $ZZ^*$ 候选。 \\
轻子对质量 & 所有 SFOS 质量 $>5$；$50<m_{Z1}<106$、$12<m_{Z2}<115$，排除低质量对并限定候选。 \\
四轻子质量、末态 & $105\leq m_{4\ell}<140$；随后要求两个电子和两个缪子，限定本项目研究范围。 \\
\bottomrule
\end{tabularx}
\end{table}
\Needspace{12\baselineskip}
\subsection{文献依据与项目差异}
\textbf{哪些条件有文献依据。} 本项目参考已发表分析的选择思路，但没有逐项复现其完整程序。
ATLAS 13 TeV 分析\cite{ATLAS13} 第 5.1--5.2 节支持接受度、
前三个动量尺度、SFOS 配对、$Z$ 质量窗口及横向碰撞参数条件的物理动机；
$12\gev$ 的次主导质量下限适用于这里的 $m_{4\ell}<140\gev$。

\textbf{项目与文献有哪些差异。} 使用上述依据时，还需保留以下区别：
\begin{itemize}[itemsep=1pt,topsep=2pt]
\item 文献使用宽松电子鉴别及详细缪子要求；本项目统一使用 tight-ID，工作点等价性未确认。
\item 本项目动量阈值允许取等号，第四轻子统一 $p_T\geq7$；文献分别要求电子 $E_T>7$、缪子 $p_T>5$，部分缪子另有要求。
\item 本项目两种隔离均 $<30\%$；文献径迹隔离 $<15\%$，量能器隔离电子 $<20\%$、缪子 $<30\%$，且有额外锥及修正规定。
\item 本项目仅取恰好四个合格轻子的 $2e2\mu$ 道；未复现文献的多候选、角分离等完整程序。该道的 $Z$ 窗口已比 $5\gev$ 条件更严格。
\end{itemize}
官方教育示例\cite{HZZCuts,HZZAnalysis}支持部分多重数、电荷、味道和动量规则，但其隔离函数调用已注释掉。
Run-1 文献\cite{ATLASRun1}支持触发匹配原则。
tight-ID 工作点及 $|z_0\sin\theta|<0.5$ mm 条件的确切外部依据仍待核实，本文将它们列为项目固定选择，
不据此声称复现已发表测量。
\FloatBarrier

\Needspace{22\baselineskip}
\section{误差计算与图中标记}\label{app:errors}

\subsection{共享事件下的保留率误差}
本节说明怎样从事件组贡献计算保留率误差。筛选后的事件也属于筛选前的样本，
所以两项产额会一起波动，计算比值误差时需计入这种联系，即协方差。

\textbf{第一步：记录每个事件组的贡献。}
对信号或背景分别计算，令第 $g$ 个事件组在筛选前的贡献为 $T_g$，在高分类别中的贡献为 $H_g$。
分别求和得到 $T=\sum_g T_g$、$H=\sum_g H_g$，保留率为 $R=H/T$。

\textbf{第二步：计算方差和协方差。}
按现有事件组计数约定，计算
\[
 V_T=\sum_g T_g^2,\qquad V_H=\sum_g H_g^2,\qquad C_{HT}=\sum_g H_gT_g.
\]
$V_T$、$V_H$ 分别用于描述筛选前后产额的方差；$C_{HT}$ 累加同一事件组在两种选择中的贡献乘积，
用于描述共享事件造成的关联。

\textbf{第三步：代入比值的误差公式。}
在 $T>0$ 时，将前两步的结果代入一阶误差传播公式：
\[
 \sigma_R^2\simeq\frac{V_H}{T^2}+\frac{H^2V_T}{T^4}-\frac{2H C_{HT}}{T^3}.
\]
前两项计入两项产额各自的方差，最后一项计入它们的协方差，不能省略。
开平方得到 $\sigma_R$，再乘以 100，得到表中以百分数展示的保留率误差。
归一化质量谱也要计入每个质量箱与总产额之间的关联，因为它们共享事件。
这里仅解释已有汇总采用的误差计算，没有为删组差值新增误差条或显著性检验。

\subsection{质量箱的统计支持标记}
阴影用来提示哪些质量箱的 MC 估计需要谨慎解读。对每个箱，分别检查筛选前和高分类别；
只要其中一种选择满足以下任一条件，就将该箱标为阴影：
\begin{enumerate}[leftmargin=1.7em,itemsep=2pt]
\item 没有事件组。
\item 带符号产额 $Y$ 或分组方差 $V$ 非正。
\item 带符号有效计数 $N_{\rm eff}=Y^2/V<20$。
\item 抵消比 $\rho=|Y|/\sum_i|u_i|<0.2$。
\end{enumerate}
其中 $u_i$ 为评估产额权重，$V$ 为该箱各事件组贡献的平方和。
计算 $N_{\rm eff}$ 是为了检查带权样本的统计支持；计算 $\rho$ 是为了检查正负权重抵消。
$\rho$ 很小时，带符号产额相对于绝对权重和很小，说明正负贡献抵消较强。
归一化谱形沿用相同标记。阴影表示上述检查条件被触发，不是产额的误差区间或置信带。

\section{实现细节与复现}\label{app:network}

\subsection{模型与汇总计算}
\textbf{训练实现。} 模型结构见第~\ref{sec:network}~节。类 \texttt{ResearchClassifier} 定义于
\nolinkurl{src/higgsml/modeling/discriminants.py}。
训练流程位于 \nolinkurl{src/higgsml/draft_v2/training.py}，超参数由冻结协议绑定。
本次网络在 CPU 上以 32 位浮点数训练。
损失使用 \texttt{binary\_cross\_entropy\_with\_logits}，直接从实数输出计算二元交叉熵，
避免先计算 sigmoid 再取对数带来的数值问题。

\Needspace{12\baselineskip}
\textbf{怎样从已有汇总复现图表。} 使用 Python 的 \texttt{json} 读取汇总文件，
NumPy 计算指标、差值和汇总，Matplotlib 绘图。按以下顺序核对和计算：
\begin{enumerate}[leftmargin=1.7em,itemsep=2pt]
\item 核对报告、冻结产物和协议的身份，确认使用的是同一组结果，并检查全部十五个模型及网络筛选前的共同产额是否完整。
\item 从 \texttt{results.json} 读取信号、背景的高低分类别产额，检查每个模型的高低产额之和是否等于网络筛选前的产额。
\item 用未舍入产额计算保留率与信背比，逐种子作删组差，再计算三个差值的中位数及最小至最大范围。
\item 最后才对展示值舍入：产额保留四位小数，百分数和百分点差保留两位小数，信背比及其差保留三位小数。
缺少种子或指标无效时明确标记，不删去失败结果后重新汇总。
\end{enumerate}
完整复现入口与构建命令见 \nolinkurl{paper/draft-v2/README.md}。

\subsection{阈值公式与教学示例}
\label{app:threshold}
本节将第~\ref{sec:threshold}~节的阈值选择写成公式，并用一个教学例子说明。
先把验证背景按分数从小到大排序。令 $w_i$ 为物理权重、$a_i=|w_i|$ 为绝对权重，下标 $(i)$ 表示排序后的位置。
计算总绝对权重的一半，再从低分开始累加，找出第一次达到该值的位置 $k$：
\[
 k=\min\left\{j:\sum_{i=1}^{j}a_{(i)}\geq\frac12\sum_i a_i\right\},
 \qquad t=s_{(k)}.
\]
使用公式时，先找到 $k$，再读取该事件的实际分数 $s_{(k)}$，作为阈值 $t$；不在相邻分数之间插值。
分数相同的事件全部按同一阈值判断。

\textbf{教学示例（非实际运行数据）。} 假设五个验证背景事件的有序分数为
$0.10,0.30,0.60,0.80,0.90$，对应绝对权重为 $1,2,1,4,2$。
总权重为 10，一半为 5。依次累加得到 $1,3,4,8,10$，第一次达到 5 时对应的分数为 0.80，因此 $t=0.80$。
按 $s\geq t$ 筛选后，分数 0.80 和 0.90 的事件进入高分类别，绝对权重占比为 $6/10=60\%$。
这个例子说明：累计权重可能直接跨过一半，加上同分事件整体处理，高分类别的权重占比不一定恰好为一半。
\FloatBarrier

\Needspace{16\baselineskip}
\section{结果来源与历史补充}\label{app:provenance}

\subsection{当前结果来源}
\textbf{使用哪一份结果。} 正文数据来自 \nolinkurl{runs/h4l-draft-v2-draft001/report/}，报告身份为
\nolinkurl{a7e610a163e5d155705014201b7ee330323e3bf69309e0e89b2496edba5d4317}，
绑定协议 \nolinkurl{config/protocols/h4l_draft_v2.json}。
\textbf{已经核对到哪一步。} 既有证据记录核对了模型与汇总的身份、文件校验、数值一致性和文稿构建，
没有重新审计原始事件载荷。详细记录见 \nolinkurl{paper/draft-v2/BASIC-20261108-EVIDENCE.md}。
文献及样本参数沿用 2026 年 10 月 30 日和 11 月 1 日的核对记录。
本轮修改解释顺序和措辞，并补充计算说明；已有公式、图表数值及图件保持原样。
没有读取事件载荷、重新训练、评分、拟合或新增外部文献核查。
CERN 发布记录将该集合限定于教育用途，不适用于科学出版，本文据此保持教育与技术研究范围。

\subsection{历史统计方法}\label{app:interval}
历史 draft001 与 validation001 曾研究区间宽度，使用的方法包括模型期望观测、剖面似然区间、配对重抽样和模型自生成伪实验。
其中曾观察到标称区间欠覆盖，尚未建立独立覆盖率校准。
本文的保留率和信背比不使用这些历史结果，这些诊断也不作为当前运行的验证。
核对历史材料时需依据产物身份，不能只看运行名称。同名运行对应的产物可以不同；原数值表、身份与限制保存在
\nolinkurl{paper/draft-v2/DRAFT001-EVIDENCE.md} 及其引用的历史材料中。

\Needspace{28\baselineskip}
\begingroup
\small
\renewcommand{\refname}{本报告参考文献}
\begin{thebibliography}{9}
\bibitem{OpenData}
ATLAS Collaboration, \emph{ATLAS 13 TeV samples collection at least four leptons
(electron or muon), for 2020 Open Data release} (2020).
\href{https://opendata.cern.ch/record/15005}{CERN Open Data record 15005},
\href{https://doi.org/10.7483/OPENDATA.ATLAS.2Y1T.TLGL}{doi:10.7483/OPENDATA.ATLAS.2Y1T.TLGL}.
两份模拟数据及发布前宽松预筛选说明的来源；使用限制见记录声明。

\bibitem{ATLAS13}
ATLAS Collaboration, \emph{Measurement of the Higgs boson coupling properties in
the $H\to ZZ^*\to4\ell$ decay channel at $\sqrt{s}=13$ TeV with the ATLAS detector},
JHEP \textbf{03} (2018) 095.
\href{https://doi.org/10.1007/JHEP03(2018)095}{doi:10.1007/JHEP03(2018)095};
\href{https://arxiv.org/pdf/1712.02304v2#page=8}{arXiv:1712.02304v2，第 5.1--5.2 节，第 8--9 页}。
附录~\ref{app:cuts} 中筛选与工作点比较的来源。

\bibitem{HZZCuts}
ATLAS Open Data outreach tools, \emph{HZZCuts.py}，版本
\nolinkurl{8ad2015be6d350060c3a183aff5802570bc9ada4}。
\href{https://github.com/atlas-outreach-data-tools/atlas-outreach-Python-uproot-framework-13tev/blob/8ad2015be6d350060c3a183aff5802570bc9ada4/HZZCuts.py#L2-L33}{固定版本源码，第 2--33 行}。
多重数、电荷、部分动量条件、味道及可选隔离定义。

\bibitem{HZZAnalysis}
ATLAS Open Data outreach tools, \emph{HZZAnalysis.py}，版本同文献~\cite{HZZCuts}。
\href{https://github.com/atlas-outreach-data-tools/atlas-outreach-Python-uproot-framework-13tev/blob/8ad2015be6d350060c3a183aff5802570bc9ada4/HZZAnalysis.py#L140-L168}{固定版本源码，第 140--168 行}。
事件权重因子、实际施加的筛选及已注释的隔离调用。

\bibitem{ATLASRun1}
ATLAS Collaboration, \emph{Measurements of Higgs boson production and couplings
in the four-lepton channel in $pp$ collisions at center-of-mass energies of
7 and 8 TeV with the ATLAS detector}, Phys.~Rev.~D \textbf{91}, 012006 (2015).
\href{https://doi.org/10.1103/PhysRevD.91.012006}{doi:10.1103/PhysRevD.91.012006};
\href{https://arxiv.org/pdf/1408.5191v3#page=7}{arXiv:1408.5191v3，第 5.1 节，印刷页码第 6--7 页}。
触发匹配原则的参考；不声称其详细阈值和工作点与本项目相同。
\bibitem{SampleInfo}
ATLAS Open Data outreach tools, \emph{infofile.py}，版本同文献~\cite{HZZCuts}。
\href{https://github.com/atlas-outreach-data-tools/atlas-outreach-Python-uproot-framework-13tev/blob/8ad2015be6d350060c3a183aff5802570bc9ada4/infofile.py}{固定版本样本参数表}。
\texttt{ggH125\_ZZ4lep} 与 \texttt{llll} 条目的 DSID、\texttt{xsec} 和 \texttt{sumw}；
不提供完整生成器或探测器模拟配置。
\end{thebibliography}
\endgroup
\end{document}
